\centerline{\bf \Huge  Правила математической драки}
\title{Правила математической драки}

1. Математическая драка — это командная игра-соревнование по решению задач. Играется командами по 4-5 человек. Каждая команда в начале игры получает список задач, которые можно решать в произвольном порядке, и первоначальный капитал в 10 тугриков.

2. Команда, считающая, что она решила задачу, пишет на листочке название команды, номер задачи и ответ на нее. После этого капитан заявляет о решении, подняв листок с номером команды и по сигналу жюри сдает заполненный бланк.

3. Первоначальная цена каждой задачи — 3 тугрика. Если команда дала неверный ответ, из ее капитала вычитается 1 тугрик, а текущая цена соответствующей задачи увеличивается на 1 тугрик. Команда, которая первой верно решила задачу, увеличивает свой капитал на цену задачи, а прием ответов по этой задаче прекращается.

4. Текущие цены задач и количество тугриков на счетах команд отмечаются на доске. Решенные задачи вычеркиваются, о снятии каждой задачи жюри также объявляет вслух.

5. Если у команды закончились тугрики, она НЕ выбывает из драки, а уходит в минус.

6. Драка заканчивается, если решены все задачи или истекло отведенное на нее время (оно объявляется в начале драки).

7. Победителем считается команда, у которой на момент окончания драки на счету больше всего тугриков (или кто меньше уйдет в минус :)). При равном числе тугриков у нескольких команд победителем считается та, которая дала больше верных ответов.